\section{Related Work}
\label{sec:related}

\paragraph{Term Rewriting}
Term rewriting~\cite{rewritesystems} has been used widely to facilitate
equational reasoning for program
optimizations~\cite{DBLP:conf/scitools/BoyleHW96,
  DBLP:journals/toplas/BrandHKO02, DBLP:conf/icfp/VisserBT98}. A term rewriting
system applies a database of semantics preserving rewrites or axioms to an input
expression to get a new expression, which may, according to some cost function,
be more profitable compared to the input. Rewrites are typically symbolic and
have a left hand side and a right hand side. To apply a rewrite to an
expression, a rewrite system implements pattern matching---if the left hand side
of a rewrite rule matches with the input expression, the system computes a
substitution which is then applied to the right-hand side of the rewrite rule.
Upon applying a rewrite rule, a rewrite system typically replaces the old
expression by the new expression. This can lead to the \textit{phase ordering}
problem--- it makes it impossible to apply a rewrite to the old expression in
the future which could have led to a more optimal result.

\paragraph{\Egraphs and E-matching}
\Egraph were originally proposed several decades ago as an
efficient data structure for maintaining congruence
closure~\cite{nelson-oppen-78, kozen-stoc77, nelson}.
\Egraphs continue to be a critical component in successful SMT
solvers where they are used for combining satisfiability theories by sharing equality
information~\cite{z3}.
 A key difference between past implementations of \egraphs and
  \egg's \egraph is our novel rebuilding algorithm that maintains
  invariants only at certain critical points~(\autoref{sec:rebuild}).
This makes \egg more efficient for the purpose of equality saturation.
  \egg implements the pattern compilation strategy introduced by de Moura et al.~\cite{ematching}
  that is used in state of the art theorem provers~\cite{z3}.
Some provers~\cite{z3, simplify} propose optimizations like mod-time,
pattern-element and inverted-path-index to find new terms and relevant patterns
for matching, and avoid redundant matches. So far, we have found \egg to be
faster than several prior \egraph implementations even without
these optimizations. They are,
however, compatible with \egg's design and could be explored in the future. Another
key difference is \egg's powerful \eclass analysis abstraction and flexible
interface. They empower the programmer to easily leverage \egraphs for problems
involving complex semantic reasoning.

\paragraph{Congruence Closure}
Our rebuilding algorithm is similar to the congruence closure algorithm
  presented by \citet{downey-cse}.
The contribution of rebuilding is not \emph{how} it restores the \egraph invariants
  but \emph{when};
  it gives the client the ability to specialize invariant restoration to a
  particular workload like equality saturation.
Their algorithm also features a worklist of merges to be processed further,
  but it is offline, i.e.,
  the algorithm processes a given set of equalities and outputs the set of
  equalities closed over congruence.
Rebuilding is adapted to the online e-graph (and equality saturation) setting,
  where rewrites frequently examine the current set of equalities and assert new ones.
Rebuilding additionally propagates e-class analysis facts (\autoref{sec:analysis}).
Despite these differences,
  the core algorithms algorithms are similar enough that theoretical results on
  offline performance characteristics \cite{downey-cse} apply to both.
We do not provide theoretical analysis of rebuilding for the online setting;
  it is likely highly workload dependent.

\paragraph{Superoptimization and Equality Saturation}
The Denali~\cite{denali} superoptimizer first demonstrated how to use \egraphs
for optimized code generation as an alternative to hand-optimized machine code
and prior exhaustive approaches~\cite{massalin}, both of which were less
scalable. The inputs to Denali are programs in a C-like language from which it
produces assembly programs. Denali supported three types of
rewrites---arithmetic, architectural, and program-specific. After applying these
rewrites till saturation, it used architectural description of the hardware to
generate constraints that were solved using a SAT solver to output a
near-optimal program. While Denali's approach was a significant improvement over
prior work, it was intended to be used on straight line code only and therefore,
did not apply to large real programs.

Equality saturation~\cite{eqsat, eqsat-llvm} developed a compiler optimization
phase that works for complex language constructs like loops and conditionals.
The first equality saturation paper used an intermediate representation called
Program Expression Graphs (PEGs) to encode loops and conditionals. PEGs have
specialized nodes that can represent infinite sequences, which allows them to
represent loops. It uses a global profitability heuristic for extraction which
is implemented using a pseudo-boolean solver. Recently, \cite{yogo-pldi20} uses
PEGs for code search.
\egg can support PEGs as a user-defined language,
and thus their technique could be ported.
%PEGs could be implemented in \egg.

%%% Local Variables:
%%% TeX-master: "egg"
%%% End:
